\section{The planner's problem}\label{sec:sp:problem}

Section~\ref{sec:sp:setup} reduced the process-level accounting to a single physical relation, the
ledger~\eqref{eq:sp:setup:ledger}, together with one inequality, $\Omega\ge0$. This section states
the planner's problem those relations imply, and derives its first-order conditions in full. The
baseline throughout is the closure in which $\Phi^I(t)$ is exogenous
(Section~\ref{sec:sp:setup:aggregatematerial}); we call it \textbf{Version B}. The alternative
closure $\Phi^I=\Omega\phi^I$, \textbf{Version A}, is taken up in Section~\ref{sec:sp:versions},
where it turns out to change the optimality conditions in one systematic and interpretable way
rather than to require a separate derivation.


\subsection{Statement}\label{sec:sp:problem:statement}

\smalltitle{Variables.} It is convenient to write the recycling technology as a constant-returns
function of treated waste and recycling capital. Let
\begin{align}
T &\equiv \varpi W,
&
R^R &= \mathcal R\left(T,K^R\right)\;\equiv\;a\!\left(\frac{K^R}{T}\right)T,
\label{eq:sp:problem:recycling-crs}
\end{align}
which is~\eqref{eq:sp:setup:recycling} rewritten. Because $a$ depends on $K^R$ and $T$ only through
their ratio, $\mathcal R$ is homogeneous of degree one, and its two partial derivatives are
\begin{align}
\mathcal R_T &= a\left(x\right)-a'\!\left(x\right)x \;\equiv\;\alpha,
&
\mathcal R_{K^R} &= a'\!\left(x\right),
&
\mathcal R &= \alpha T+a'K^R .
\label{eq:sp:problem:recycling-derivatives}
\end{align}
We refer to $\alpha=a\left(1-\varepsilon\right)$ as the \emph{marginal recycling yield}: the extra
tonne of secondary material obtained from one extra tonne of treated waste, holding recycling
capital fixed. It is strictly below the average yield $a$ whenever $\varepsilon>0$, and the
assumption $\varepsilon<1$ in~\eqref{eq:sp:setup:recycling} is exactly the requirement
$\alpha>0$.

The planner chooses paths for the seven controls
\begin{align}
\left\{C,\;I,\;N,\;D,\;\varpi,\;K^R,\;W\right\},
\label{eq:sp:problem:controls}
\end{align}
which determine the remaining flows through the definitions
\begin{align}
R^R &= \mathcal R\left(\varpi W,K^R\right),
&
R &= N+R^R,
&
Y &= F\left(K-K^R,\,R,\,P\right),
&
G &= I+\delta K,
\label{eq:sp:problem:definitions}
\end{align}
and it carries the five states
\begin{align}
\left\{K,\;S,\;X,\;P,\;M^K\right\},
\label{eq:sp:problem:states}
\end{align}
with $K_0,S_0,X_0,P_0,M^K_0$ given.

\smalltitle{Why $W$ is listed as a control.} The ledger determines $W$ residually, so it is not an
independent object; listing it among the controls and imposing the ledger as an equality constraint
is the equivalent formulation that keeps the shadow price of waste visible. Substituting the solved
ledger~\eqref{eq:sp:setup:ledger-solved} and dropping $W$ from~\eqref{eq:sp:problem:controls} yields
the same allocation, but the multiplier $p^W$ --- the object every environmental instrument in
Section~\ref{sec:sp:versions} and in the market economy is built from --- would then have to be
recovered afterwards. Remark~\ref{rem:sp:reduced} states the equivalence.

\begin{problem}[Planner, Version B]\label{prob:sp:planner}
Choose $\left\{C,I,N,D,\varpi,K^R,W\right\}_{t\ge0}$ to maximise
$U_0=\int_0^\infty\left[u(C)-v(P)\right]e^{-\rho t}dt$ subject to, at every $t$,
\begin{align}
&\text{\emph{goods:}}
&
F\left(K-K^R,R,P\right) &= C+I+\delta K+C^N(N,S)+C^D(D,X)+c^W\!\left(\varpi\right)W,
\label{eq:sp:problem:goods}
\\[0.2em]
&\text{\emph{ledger:}}
&
W &= R-\Phi^I G+\delta M^K+\mathcal N,
\label{eq:sp:problem:ledger}
\\[0.2em]
&\text{\emph{mass:}}
&
\Phi^I G &\le R,
\label{eq:sp:problem:mass}
\end{align}
and to the laws of motion
\begin{align}
\dot K &= I,
&
\dot S &= D-N,
&
\dot X &= D,
&
\dot P &= \Xi-\theta(P)P,
&
\dot M^K &= \Phi^I G-\delta M^K,
\label{eq:sp:problem:states-motion}
\end{align}
with $\mathcal N=\Omega^{N,S}N+\Omega^{D,S}D$, $c^W(\varpi)=c^c+c^T(\varpi)$, the definitions
in~\eqref{eq:sp:problem:definitions}, the emission flow
\begin{align}
\Xi &= \left(1-\varpi+d^W\varpi\right)W-d^WR^R,
\label{eq:sp:problem:Xi}
\end{align}
and $C,I+\delta K,N,D,K^R,W\ge0$, $\varpi\in\left[0,1\right]$.
\end{problem}

Three features of this statement deserve marking before any algebra.

First, \eqref{eq:sp:problem:ledger} is an \emph{identity}, not a technology. It says where the mass
went, and it holds along every conceivable path; what makes it bite is that it ties the waste flow
$W$ --- which is costly to collect, costly to treat and damaging when released --- to the material
flow $R$ that production requires. Nothing in it can be relaxed by choosing differently; it can only
be made less costly by choosing a path on which less mass has to be moved.

Second, \eqref{eq:sp:problem:mass} is the only part of the materials accounting that restricts the
feasible set, and it does so as an inequality. This is the point on which the present formulation
departs most sharply from the reduced-form treatment: imposing the materials balance as an
\emph{equality} constraint with a freely signed multiplier would attach a nonzero shadow price to
mass conservation by construction, and that price would show up as a wedge on output and on
consumption. Here the multiplier on~\eqref{eq:sp:problem:mass} is zero except at the corner where
every tonne entering production is embodied in new capital.

Third, $\Xi$ is written in the form~\eqref{eq:sp:problem:Xi}, linear in $W$ and $R^R$ separately, so
that no derivative of $a(\cdot)$ ever appears inside a pollution term. All the curvature of the
recycling technology enters through $\mathcal R$ and through $\alpha$ and $a'$ alone.


\subsection{The Hamiltonian}\label{sec:sp:problem:hamiltonian}

Write the current-value Hamiltonian with the static constraints adjoined. It is convenient to
normalise every multiplier by $\lambda$, the multiplier on the goods
constraint~\eqref{eq:sp:problem:goods}, so that all shadow prices below are measured in units of the
final good rather than in utils. Accordingly, let the current-value costates be
\begin{align}
\underbrace{\lambda q}_{K},
&&
\underbrace{\lambda p^S}_{S},
&&
\underbrace{\lambda p^X}_{X},
&&
\underbrace{-\lambda p^P}_{P},
&&
\underbrace{-\lambda p^M}_{M^K},
\label{eq:sp:problem:costates}
\end{align}
the two minus signs anticipating that both $P$ and $M^K$ are liabilities, so that $p^P$ and $p^M$
read as shadow \emph{costs}; the signs are verified in
Section~\ref{sec:sp:solution:signs}, not assumed. Let $\lambda p^W$ be the multiplier on the
ledger~\eqref{eq:sp:problem:ledger} written with $W$ on the left, and $\lambda\eta\ge0$ the
multiplier on~\eqref{eq:sp:problem:mass}. Then
\begin{align}
\mathcal L
={}& u(C)-v(P)
\notag\\
&+\lambda\Big[F\left(K-K^R,R,P\right)-C-I-\delta K-C^N(N,S)-C^D(D,X)-c^W\!\left(\varpi\right)W\Big]
\notag\\
&+\lambda q\,I
\;+\;\lambda p^S\left(D-N\right)
\;+\;\lambda p^X D
\;-\;\lambda p^P\left[\Xi-\theta(P)P\right]
\;-\;\lambda p^M\left[\Phi^IG-\delta M^K\right]
\notag\\
&+\lambda p^W\Big[W-N-R^R+\Phi^IG-\delta M^K-\mathcal N\Big]
\;+\;\lambda\eta\Big[N+R^R-\Phi^IG\Big],
\label{eq:sp:problem:lagrangian}
\end{align}
with $R=N+R^R$, $R^R=\mathcal R\left(\varpi W,K^R\right)$ and $G=I+\delta K$ substituted throughout.

The sign convention on $p^W$ is worth stating explicitly, because it is the price everything else
is written in terms of. Perturb the ledger to $W=R-\Phi^IG+\delta M^K+\mathcal N-\kappa$, that is,
suppose $\kappa$ tonnes of waste could be made to disappear at no cost. Then
$\partial\mathcal L^*/\partial\kappa=\lambda p^W$, so $p^W>0$ says exactly that waste is socially
costly, and $p^W$ measures how costly, per tonne, in final-good units.


\subsection{First-order conditions}\label{sec:sp:problem:foc}

Differentiating~\eqref{eq:sp:problem:lagrangian} with respect to each of the seven controls, and
dividing throughout by $\lambda$, gives the following system. Two composite objects recur often
enough to be named at the outset:
\begin{align}
B &\equiv F_R-p^W+d^Wp^P+\eta,
&
p^N &\equiv C^N_N+p^S+\Omega^{N,S}p^W .
\label{eq:sp:problem:B}
\end{align}
$B$ is the net social value of one tonne of \emph{secondary} material: what it produces, less the
waste it will itself become, plus the emission its recovery avoids, plus the mass constraint it
relaxes. $p^N$ is the full marginal social cost of one tonne of \emph{virgin} material delivered to
production: the resources spent extracting it, the rent on the reserve it depletes, and the waste
cost of the overburden its extraction displaces. Section~\ref{sec:sp:solution} argues that the
relation between these two objects is the single most informative equation in the model.

\smalltitle{Consumption.}
\begin{align}
u'\!\left(C\right) &= \lambda .
\label{eq:sp:foc:C}
\end{align}
No wedge: at the optimum the marginal utility of consumption equals the shadow value of a unit of
the final good, exactly as in an economy with no materials at all.

\smalltitle{Investment.}
\begin{align}
q &= 1+\Phi^I\left(p^M-p^W+\eta\right).
\label{eq:sp:foc:I}
\end{align}
The shadow price of installed capital is the goods it costs, \emph{less} the waste that embodying
$\Phi^I$ tonnes in it removes from today's waste stream, \emph{plus} the discounted waste those same
tonnes will become when the capital depreciates. Capital is a temporary materials sink, and $q$
prices the net value of that service.

\smalltitle{Extraction.}
\begin{align}
F_R+\eta-p^W &= C^N_N+p^S+\Omega^{N,S}p^W,
\qquad\text{equivalently}\qquad
B &= p^N+d^Wp^P .
\label{eq:sp:foc:N}
\end{align}

\smalltitle{Exploration.}
\begin{align}
p^S+p^X &= C^D_D+\Omega^{D,S}p^W .
\label{eq:sp:foc:D}
\end{align}

\smalltitle{Treatment share.}
\begin{align}
\alpha B+\left(1-d^W\right)p^P &= \frac{dc^T}{d\varpi},
\label{eq:sp:foc:varpi}
\end{align}
for interior $\varpi\in\left(0,1\right)$, with the usual inequalities at the corners.

\smalltitle{Recycling capital.}
\begin{align}
F_K &= a'\!\left(x\right)B .
\label{eq:sp:foc:KR}
\end{align}

\smalltitle{Waste.}
\begin{align}
p^W &= c^W\!\left(\varpi\right)+\left(1-\varpi+d^W\varpi\right)p^P-\varpi\alpha B .
\label{eq:sp:foc:W}
\end{align}

\smalltitle{Complementary slackness.}
\begin{align}
\eta &\ge0,
&
\Phi^IG &\le R,
&
\eta\left[R-\Phi^IG\right] &= 0 .
\label{eq:sp:foc:cs}
\end{align}

Two of these are worth expanding on immediately, because their form is not obvious.

\eqref{eq:sp:foc:W} defines $p^W$ implicitly, since $B$ contains $p^W$. Solving out, the shadow
cost of a tonne of gross waste satisfies
\begin{align}
p^W\left(1+\Omega^{N,S}\varpi\alpha\right)
&= c^W\!\left(\varpi\right)+\left(1-\varpi+d^W\varpi\right)p^P
-\varpi\alpha\left(C^N_N+p^S+d^Wp^P\right),
\label{eq:sp:foc:W-solved}
\end{align}
where~\eqref{eq:sp:foc:N} was used to replace $B$. A tonne of waste costs what it costs to collect
and treat, plus the damage from the part of it that reaches the environment, minus the virgin
extraction it displaces at the margin. \emph{Nothing in~\eqref{eq:sp:foc:W-solved} involves any
$\omega^j$ or $\Omega$}: the shadow price of waste is a function of extraction costs, resource
rents, damages and the recycling technology alone.

\eqref{eq:sp:foc:varpi} is the treatment margin. Raising $\varpi$ by one point costs
$dc^T/d\varpi$ per tonne of waste and buys two things: $\alpha$ tonnes of secondary material worth
$B$ each, and a reduction of the emitted fraction from $1$ to $d^W$ on the tonne now treated. Since
$c^T(0)=\left.dc^T/d\varpi\right|_{0}=0$ by~\eqref{eq:sp:setup:waste-cost}, the right-hand side
vanishes at $\varpi=0$ while the left-hand side is strictly positive whenever $p^P>0$; so
\emph{some} treatment is always optimal, and the $\varpi=0$ corner is never reached. The upper
corner $\varpi=1$ is reached iff $\alpha B+\left(1-d^W\right)p^P\ge\left.dc^T/d\varpi\right|_{1}$.


\subsection{Costate equations}\label{sec:sp:problem:costates}

Each costate obeys $\dot\nu=\rho\nu-\partial\mathcal L/\partial(\text{state})$ with $\nu$ the
current-value costate in~\eqref{eq:sp:problem:costates}. Because the costates carry the factor
$\lambda$, each such equation produces a term $\dot\lambda/\lambda$, and it is worth defining the
resulting discount rate once:
\begin{align}
r &\equiv \rho-\frac{\dot\lambda}{\lambda}.
\label{eq:sp:problem:r}
\end{align}
Since $\lambda=u'(C)$ by~\eqref{eq:sp:foc:C}, $r$ is the familiar consumption rate of interest,
$r=\rho+\sigma(C)\dot C/C$ with $\sigma\equiv-Cu''/u'$; and because every shadow price has been
normalised by $\lambda$, \emph{every} price below discounts at this same $r$. That uniformity is not
automatic --- it is a consequence of the normalisation, and it fails in Version A, where
$\lambda\neq u'(C)$.

\smalltitle{Capital.} Using $\partial\mathcal L/\partial K=\lambda\left[F_K-\delta-\delta\Phi^I
\left(p^M-p^W+\eta\right)\right]$ and~\eqref{eq:sp:foc:I} to collapse the bracket to
$F_K-\delta q$,
\begin{align}
r &= \frac{F_K+\dot q}{q}-\delta .
\label{eq:sp:costate:K}
\end{align}
Depreciation is charged at $q$, not at $1$: what wears out is a unit of installed capital, whose
social value is $q$, not a unit of the final good. Together with~\eqref{eq:sp:problem:r} this is the
consumption Euler equation,
\begin{align}
\sigma\left(C\right)\frac{\dot C}{C} &= \frac{F_K+\dot q}{q}-\delta-\rho .
\label{eq:sp:costate:euler}
\end{align}

\smalltitle{Reserves.} With $\partial\mathcal L/\partial S=-\lambda C^N_S$,
\begin{align}
\dot p^S &= r\,p^S+C^N_S,
\qquad\text{so}\qquad
\frac{\dot p^S}{p^S} = r+\frac{C^N_S}{p^S}<r .
\label{eq:sp:costate:S}
\end{align}
This is the modified Hotelling rule: the in-situ rent grows more slowly than the interest rate,
by exactly the marginal cost saving $-C^N_S>0$ that a larger known reserve confers on extraction.

\smalltitle{Discoveries.} With $\partial\mathcal L/\partial X=-\lambda C^D_X$,
\begin{align}
\dot p^X &= r\,p^X+C^D_X,
\qquad\text{so}\qquad
p^X\left(t\right) = -\int_t^\infty C^D_X\!\left(s\right)e^{-\int_t^s r\,du}\,ds<0 .
\label{eq:sp:costate:X}
\end{align}
Cumulative discovery is a bad: each unit found today raises the cost of finding the next one
forever after, and $-p^X$ is the present value of that penalty.

\smalltitle{Pollution.} Write the marginal damage of the stock, in final-good units, as
\begin{align}
d &\equiv \frac{v'\!\left(P\right)}{u'\!\left(C\right)}-F_P\;>\;0,
\label{eq:sp:problem:damage}
\end{align}
the sum of the direct disutility and the output loss. Then
\begin{align}
\dot p^P &= \Big[r+\theta\!\left(P\right)+\theta'\!\left(P\right)P\Big]p^P-d,
\qquad\text{so}\qquad
p^P\left(t\right)=\int_t^\infty d\left(s\right)e^{-\int_t^s\left[r+\theta+\theta'P\right]du}\,ds .
\label{eq:sp:costate:P}
\end{align}
The effective discount rate on damages is the interest rate plus the marginal decay rate
$\partial\left[\theta(P)P\right]/\partial P=\theta+\theta'P$. Because $\theta'<0$
by~\eqref{eq:sp:setup:pollution-motion}, sink saturation lowers that rate below the average decay
rate $\theta$, and $p^P$ is correspondingly larger than a naive $d/\left(r+\theta\right)$ calculation
would suggest. Stability of the pollution block requires $\theta+\theta'P>0$; this is an assumption
on $\theta$, and it is maintained throughout.

\smalltitle{Embodied material.} With
$\partial\mathcal L/\partial M^K=\lambda\delta\left(p^M-p^W\right)$,
\begin{align}
\dot p^M &= \left(r+\delta\right)p^M-\delta p^W,
\qquad\text{so}\qquad
p^M\left(t\right)=\int_t^\infty \delta p^W\!\left(s\right)e^{-\left(\int_t^s r\,du\right)-\delta\left(s-t\right)}ds .
\label{eq:sp:costate:M}
\end{align}
$p^M$ is the present value of the waste that the currently embodied stock will release as it
depreciates, discounted at $r+\delta$ because a fraction $\delta$ of the liability is discharged
each instant. It is positive whenever $p^W$ is, confirming the sign anticipated
in~\eqref{eq:sp:problem:costates}, and it is strictly below the contemporaneous $p^W$ whenever
$r>0$ --- postponing a tonne of waste is worth something.

\smalltitle{Transversality.} The problem is closed by
\begin{align}
\lim_{t\to\infty}e^{-\rho t}\lambda q K &= 0,
&
\lim_{t\to\infty}e^{-\rho t}\lambda p^S S &= 0,
&
\lim_{t\to\infty}e^{-\rho t}\lambda p^X X &= 0,
\notag\\
\lim_{t\to\infty}e^{-\rho t}\lambda p^P P &= 0,
&
\lim_{t\to\infty}e^{-\rho t}\lambda p^M M^K &= 0 .
&&
\label{eq:sp:problem:tvc}
\end{align}

\begin{remark}[The reduced formulation]\label{rem:sp:reduced}
Substituting the solved ledger~\eqref{eq:sp:setup:ledger-solved} into the goods constraint and
dropping $W$ from the control set gives a problem in six controls and the same five states, whose
first-order conditions are~\eqref{eq:sp:foc:C}--\eqref{eq:sp:foc:KR} with $p^W$ replaced everywhere
by the right-hand side of~\eqref{eq:sp:foc:W}. The two formulations are equivalent because
\eqref{eq:sp:foc:W} is precisely the condition that makes the substitution valid; what is lost is
the price itself. Since $p^W$ is the object the market economy implements through a waste tax, and
since the sign of $p^W$ turns out to determine the sign of two separate instruments simultaneously
(Section~\ref{sec:sp:solution:signs}), it is worth keeping.
\end{remark}

\begin{remark}[What is not verified]\label{rem:sp:soc}
The conditions above are necessary. Sufficiency requires concavity of the maximised Hamiltonian in
the states, which is not established here: $F$ is concave and $u,-v$ are, but $\mathcal R$ enters
through a constant-returns function of two controls, and the ledger is bilinear in $\varpi$ and $W$.
Uniqueness of the optimum is likewise not established, and it is invoked later when the market
implementation is claimed to reproduce \emph{the} planner's allocation. Both are open.
\end{remark}
